%%%%%%%%%%% Conclusion %%%%%%%%%%%%%


Au terme de notre alternance aux Arts décoratifs (MAD), nous avions pour objectif d'apporter une solution permettant de pallier au problème du \enquote{vrac numérique}, qui représente l'une des principales préoccupations des archivistes et les autres utilisateurs du serveur. L'accumulation des informations dans l'espace de stockage affecte directement l’intelligibilité et la pérennité du patrimoine institutionnel. On remarque donc que \enquote{les défis posés par la gestion de grandes quantités de données numériques exigent de nouvelles approches computationnelles permettant le traitement et l'accès aux données}\footnote{\cite{saad2026intelligence} p. 83.}. En outre, seule l'utilisation de la méthode traditionnelle de traitement archivistique ne permet pas de répondre aux enjeux liés à la résolution de ce phénomène : c'est ce que notre recherche a cherché à démontrer. Nous interrogeons la manière dont l'expertise archivistique, l'audit technique de l'espace de stockage et le recours au traitement algorithmique se croisent. Notre travail met en lumière l'alliance entre ces trois dimensions pour redonner du sens à la masse que représente le vrac.\\[0.1cm]

La mise en place d'un système d'automatisation des tâches avec scripts Python nous aide à décompresser, trier et ranger les fichiers dans le plan de classement de l'institution. Nous avons donc conçu cet outil, pour qu'il reflète nos pratiques documentaires. Son but est de montrer qu'il est tout à fait possible pour l'archiviste de se confronter aux défis du traitement massif des archives numériques alliés aux défis du traitement automatisé. Il faut comprendre par là que le programme informatique lui seul ne dispose pas d'un intérêt patrimonial s'il n'est pas guidé par la réflexion de l'archiviste. Il contribue à la transformation d'un espace de stockage en un fonds d'archives structuré et exploitable. Cela témoigne de l'évolution du métier l'archiviste au XVIe siècle, désormais confronté à problématique liées à la fois au traitement et à l'accès à l'information.\\[0.1cm]

Notre recherche a pour but de montrer l'évolution profonde du métier de l'archiviste et de ses pratiques. Toutefois, même face au recours des nouveaux outils de traitement technologique, ceux-ci ne le remplacent pas. On peut penser qu'en automatisant ses tâches on n'a plus forcément besoin de l'humain. En pratique, c'est le contraire parce que les règles de conformités sont entièrement rédigées par l'humain. L'archiviste effectue le contrôle qualité pour s'assurer de sa bonne exécution.\\[0.1cm]

Pour démontrer de façon pratique le rapport de complémentarité des trois notions évoquées, nous sommes passés par un traitement en triplet : partant de l'audit technique des dossiers exposition, en passant par la conception des scripts Python et tout ceci encadré par notre réflexion archivistique. Nous avons appliqué cette approche à un corpus réel de dossiers d'exposition. Ce qui nous a permis de nous assurer de l’exactitude du respect des règles prédéfinies par l'humain au programme, mais aussi d'ajuster et de valider notre réflexion. De plus, le fait pour d'avoir fait des erreurs nous a permis de rendre le traitement plus robuste et ciblé. 

Pour conclure, notre propos, nous pouvons dire que cette expérimentation au MAD va au-delà d'un simple exercice technique. Elle constitue le moyen de confronter la théorie archivistique à la réalité opérationnelle des grands espaces de stockage : il est donc tout à fait possible de concevoir des outils sur-mesure adaptés à des cas particuliers. Le prototype mis en place permet à l'institution de mieux préparer les données avant leur transfert sous la forme de SIP vers le système d'archivage dédié. Il peut être appliqué à tous les cas d'études, pourvus qu'ils soient confrontés aux problèmes de masse documentaire mal organisée.






